% GENERATED by docs/manual/build.py - do not edit; change the source it names instead.
% source: scripts/watchover.py (SKILL_COMMANDS, TOOLS and each tool's argparse parser)
\chapter{Skill commands (talking to your AI agent)}\label{ch:skill-commands}
These commands are typed in a conversation with any AI agent that has the skill, the MCP prompts
or the prompt pack (Chapter~\ref{doc-agents}): write \texttt{akeyless-watchover <command> [args]} or your
agent's skill syntax (\texttt{/akeyless-watchover ...}, \texttt{\$akeyless-watchover ...}); MCP clients list them
as prompts \texttt{watchover-<command>}. You can also simply describe the task -- the skill also triggers
without a command. Arguments are free text.
Every section ends with an \emph{example session}: the prompt and an abridged version of the
answer. Answers depend on the agent, the model and the conversation, so the sessions are illustrative; the
artefacts they produce are the files and commands documented in Part~\ref{part:reference}.

\begin{center}\small
\begin{tabularx}{\linewidth}{@{}l>{\raggedright\arraybackslash}X@{}}
\toprule
\textbf{Command} & \textbf{Purpose} \\ \midrule
\hyperref[skill:help]{\texttt{help}} & Show this guide, or the details of one skill command. \\
\hyperref[skill:design]{\texttt{design}} & Translate an organisation/structure diagram into Akeyless configuration. \\
\hyperref[skill:change]{\texttt{change}} & Change existing access with the smallest possible model diff. \\
\hyperref[skill:review]{\texttt{review}} & Review a proposed change for least privilege, guard violations and blast radius. \\
\hyperref[skill:pipeline]{\texttt{pipeline}} & Set up the governance repository: bootstrap identity, CI, CODEOWNERS, state. \\
\hyperref[skill:import]{\texttt{import}} & Export existing config (RBAC, auth methods incl. OIDC, folders, custom fields, key policies, PKI issuers, account settings) and translate it into Terraform. \\
\hyperref[skill:conjur]{\texttt{conjur}} & Migrate from CyberArk Conjur: secrets via the Gateway, RBAC via translation. \\
\hyperref[skill:pki]{\texttt{pki}} & Design private/public CA issuers and who may request certificates. \\
\hyperref[skill:harden]{\texttt{harden}} & Account-level hardening. \\
\hyperref[skill:gateway]{\texttt{gateway}} & Gateway admin plane and audit-log forwarding. \\
\hyperref[skill:verify]{\texttt{verify}} & Prove who can access what. \\
\hyperref[skill:cli]{\texttt{cli}} & Exact Akeyless CLI commands for a task (break-glass, bootstrap, one-off checks). \\
\hyperref[skill:explain]{\texttt{explain}} & Explain an Akeyless concept precisely (sub-claims, deny, wildcards, anti-escalation...). \\
\bottomrule
\end{tabularx}
\end{center}

\section{\texttt{help}}\label{skill:help}
\begin{description}[style=nextline,leftmargin=2em,font=\normalfont\bfseries]
\item[Syntax] \texttt{/\allowbreak{}akeyless-\allowbreak{}watchover help [command]}
\item[Arguments] [command]
\item[Purpose] Show this guide, or the details of one skill command.
\item[Delivers] The command table, parameters and examples.
\item[Example] \texttt{/\allowbreak{}akeyless-\allowbreak{}watchover help design}
\end{description}
\IfFileExists{sessions/help.tex}{% Illustrative example (hand-written); answers vary with the agent and model.
\begin{session}{Example}
\you akeyless-watchover help change

\agent Shows the \texttt{change} entry: arguments (\emph{an access request in plain words}), what it
delivers (a minimal diff of \texttt{model/roles/<bu>.yaml}, new expectations, who must approve) and
an example. Without an argument: the table of all skill commands and the kit commands.
\end{session}
}{}

\section{\texttt{design}}\label{skill:design}
\begin{description}[style=nextline,leftmargin=2em,font=\normalfont\bfseries]
\item[Syntax] \texttt{/\allowbreak{}akeyless-\allowbreak{}watchover design \textless{}diagram \textbar{}\allowbreak{} description of BUs,\allowbreak{} apps,\allowbreak{} environments,\allowbreak{} personas\textgreater{}}
\item[Arguments] \textless{}diagram \textbar{} description of BUs, apps, environments, personas\textgreater{}
\item[Purpose] Translate an organisation/structure diagram into Akeyless configuration.
\item[Delivers] Path tree, access matrix, YAML model (roles, auth methods, folders), access expectations, plan expectations, CLI equivalent.
\item[Example] \texttt{/\allowbreak{}akeyless-\allowbreak{}watchover design \textless{}paste or attach the diagram\textgreater{}}
\end{description}
\IfFileExists{sessions/design.tex}{% Illustrative example session (hand-written). See Section \ref{sec:how-examples}.
\begin{session}{Example session}
\you akeyless-watchover design \emph{[attached diagram: BU Payments with apps checkout, refunds; prod and nonprod; Entra groups; EKS workloads]}

\agent \textbf{1. Path tree}\quad \texttt{/payments/\{checkout,refunds\}/\{prod,nonprod\}/...},
\texttt{/payments/<app>/prod/break-glass/...}

\textbf{2. Access matrix}\quad payments-admin: nonprod CRUD, prod read; checkout-prod-pod: own prod
read; auditor: list everything; everyone else: nothing.

\textbf{3. Model}\quad \texttt{model/roles/payments.yaml} (\texttt{team-admin}, \texttt{checkout-prod-workload}, \ldots),
\texttt{model/auth\_methods.yaml} (OIDC with \texttt{force\_sub\_claims}, K8s bound to namespace
\texttt{payments}), \texttt{model/folders.yaml}.

\textbf{4. Expectations}\quad 11 checks including the negative ones (break-glass, other BUs).

\textbf{5. Plan expectations and CLI}\quad ``+ 2 roles, + 2 auth methods, + 3 associations, no
replacements''; the equivalent \texttt{akeyless create-role / set-role-rule / assoc-role-am} commands.
\end{session}
}{}

\section{\texttt{change}}\label{skill:change}
\begin{description}[style=nextline,leftmargin=2em,font=\normalfont\bfseries]
\item[Syntax] \texttt{/\allowbreak{}akeyless-\allowbreak{}watchover change \textless{}access request in plain words\textgreater{}}
\item[Arguments] \textless{}access request in plain words\textgreater{}
\item[Purpose] Change existing access with the smallest possible model diff.
\item[Delivers] Diff of model/roles/\textless{}bu\textgreater{}.yaml (and auth\_methods.yaml if needed), new expectations, who must approve (CODEOWNERS).
\item[Example] \texttt{/\allowbreak{}akeyless-\allowbreak{}watchover change give fraud-\allowbreak{}analysts read on /\allowbreak{}payments/\allowbreak{}+/\allowbreak{}prod/\allowbreak{}*}
\end{description}
\IfFileExists{sessions/change.tex}{% Illustrative example session (hand-written). See Section \ref{sec:how-examples}.
\begin{session}{Example session}
\you /akeyless-watchover change give the fraud-analysts Entra group read on refunds prod

\claude Smallest diff -- a new role, because no existing role has the right audience:
\begin{lstlisting}
+  /governance/roles/payments/fraud-analysts:
+    description: Fraud analysts read refunds prod
+    rules:
+      - { path: "/payments/refunds/prod/*", capabilities: [read, list] }
+      - { path: "/payments/refunds/prod/break-glass/*", capabilities: [deny] }
+    associations:
+      - auth_method: /governance/auth/sso-entra-oidc
+        sub_claims: { groups: [grp-fraud-analysts] }
\end{lstlisting}
New expectations (one positive, two negative), the expected plan (+1 role, +1 association) and the
approvers from CODEOWNERS: payments owners and the platform team.
\end{session}
}{}

\section{\texttt{review}}\label{skill:review}
\begin{description}[style=nextline,leftmargin=2em,font=\normalfont\bfseries]
\item[Syntax] \texttt{/\allowbreak{}akeyless-\allowbreak{}watchover review \textless{}model files,\allowbreak{} merge-\allowbreak{}request diff or plan output\textgreater{}}
\item[Arguments] \textless{}model files, merge-request diff or plan output\textgreater{}
\item[Purpose] Review a proposed change for least privilege, guard violations and blast radius.
\item[Delivers] Findings by severity, failing guards, missing expectations, plan red flags (replacements).
\item[Example] \texttt{/\allowbreak{}akeyless-\allowbreak{}watchover review \textless{}paste the plan\textgreater{}}
\end{description}
\IfFileExists{sessions/review.tex}{% Illustrative example session (hand-written). See Section \ref{sec:how-examples}.
\begin{session}{Example session}
\you akeyless-watchover review \emph{[pasted plan]}

\agent \textbf{High}: \texttt{akeyless\_auth\_method\_k8s.this["/governance/auth/k8s-prod"]} \emph{must be
replaced} -- the name changed; consumers lose their access ID. Revert the rename or plan a cut-over.
\textbf{Medium}: rule \texttt{/payments/*} with \texttt{update} on \texttt{team-admin} widens access
to every app's prod; no expectation covers it. \textbf{OK}: 2 in-place rule updates match the request.
\end{session}
}{}

\section{\texttt{pipeline}}\label{skill:pipeline}
\begin{description}[style=nextline,leftmargin=2em,font=\normalfont\bfseries]
\item[Syntax] \texttt{/\allowbreak{}akeyless-\allowbreak{}watchover pipeline [github \textbar{}\allowbreak{} gitlab] [repository]}
\item[Arguments] [github \textbar{} gitlab] [repository]
\item[Purpose] Set up the governance repository: bootstrap identity, CI, CODEOWNERS, state.
\item[Delivers] bootstrap.tfvars, CI file, variables to create, step-by-step order.
\item[Example] \texttt{/\allowbreak{}akeyless-\allowbreak{}watchover pipeline gitlab group/\allowbreak{}akeyless-\allowbreak{}governance}
\end{description}
\IfFileExists{sessions/pipeline.tex}{% Illustrative example (hand-written); answers vary with the agent and model.
\begin{session}{Example}
\you akeyless-watchover pipeline gitlab acme/platform/akeyless-governance

\agent \texttt{bootstrap.tfvars} (\texttt{platform = "gitlab"}, \texttt{repository = "acme/platform/akeyless-governance"},
\texttt{gitlab\_url}), the order of steps (bootstrap as admin $\rightarrow$ CI variables $\rightarrow$
copy \texttt{ci/gitlab/.gitlab-ci.yml} $\rightarrow$ protected branch and environment $\rightarrow$
CODEOWNERS), the list of CI variables to create and which ones must be masked.
\end{session}
}{}

\section{\texttt{import}}\label{skill:import}
\begin{description}[style=nextline,leftmargin=2em,font=\normalfont\bfseries]
\item[Syntax] \texttt{/\allowbreak{}akeyless-\allowbreak{}watchover import [what to export: everything /\allowbreak{} a type /\allowbreak{} one user's view] [export directory or REPORT.md]}
\item[Arguments] [what to export: everything / a type / one user's view] [export directory or REPORT.md]
\item[Purpose] Export existing config (RBAC, auth methods incl. OIDC, folders, custom fields, key policies, PKI issuers, account settings) and translate it into Terraform.
\item[Delivers] The exact \texttt{watchover export} command (model or -{}-hcl mode, filters), then resolution of every REPORT.md item and the import checklist until the plan shows imports only.
\item[Example] \texttt{/\allowbreak{}akeyless-\allowbreak{}watchover import the OIDC auth method and everything alice@example.com can reach}
\end{description}
\IfFileExists{sessions/import.tex}{% Illustrative example session (hand-written). See Section \ref{sec:how-examples}.
\begin{session}{Example session}
\you /akeyless-watchover import the OIDC auth method and everything alice@example.com can reach

\claude The command to run:
\begin{lstlisting}
scripts/watchover export --out export/ --profile admin --types roles,auth-methods \
  --persona-auth-method /governance/auth/sso-entra-oidc \
  --persona-claim email=alice@example.com --persona-claim groups=grp-akl-payments-admins
\end{lstlisting}
Then, given \texttt{export/REPORT.md}: one decision per report item (client secret into
\texttt{TF\_VAR\_auth\_method\_secrets}, a temporary rule that will not be imported, an
association to an unmanaged auth method), where to move each file, and ``iterate until the plan
shows imports only''.
\end{session}
}{}

\section{\texttt{conjur}}\label{skill:conjur}
\begin{description}[style=nextline,leftmargin=2em,font=\normalfont\bfseries]
\item[Syntax] \texttt{/\allowbreak{}akeyless-\allowbreak{}watchover conjur \textless{}Conjur policy YAML\textgreater{}}
\item[Arguments] \textless{}Conjur policy YAML\textgreater{}
\item[Purpose] Migrate from CyberArk Conjur: secrets via the Gateway, RBAC via translation.
\item[Delivers] Translated roles, bindings to fill (auth method + sub-claims), migration resource, cut-over plan.
\item[Example] \texttt{/\allowbreak{}akeyless-\allowbreak{}watchover conjur \textless{}attach policy files\textgreater{}}
\end{description}
\IfFileExists{sessions/conjur.tex}{% Illustrative example (hand-written); answers vary with the agent and model.
\begin{session}{Example}
\you akeyless-watchover conjur \emph{[attached root.yml and prod.yml]}

\agent Runs (or tells you to run) \texttt{scripts/watchover conjur root.yml prod.yml --out export/conjur},
explains each translated role (Conjur layer/group $\rightarrow$ Akeyless role on the migrated
paths), and for every entry of \texttt{bindings.yaml} proposes the auth method and sub-claims
(e.g.\ host \texttt{prod/payments/checkout} $\rightarrow$ K8s auth, namespace \texttt{payments}).
Plus the Gateway migration entry for \texttt{model/gateway.yaml} and a cut-over plan.
\end{session}
}{}

\section{\texttt{pki}}\label{skill:pki}
\begin{description}[style=nextline,leftmargin=2em,font=\normalfont\bfseries]
\item[Syntax] \texttt{/\allowbreak{}akeyless-\allowbreak{}watchover pki \textless{}certificate requirements\textgreater{}}
\item[Arguments] \textless{}certificate requirements\textgreater{}
\item[Purpose] Design private/public CA issuers and who may request certificates.
\item[Delivers] pki.yaml (CA keys, issuers), requester roles, expiry forwarder, CLI to issue/revoke.
\item[Example] \texttt{/\allowbreak{}akeyless-\allowbreak{}watchover pki internal TLS for *.payments.internal,\allowbreak{} 90 days,\allowbreak{} auto-\allowbreak{}renew}
\end{description}
\IfFileExists{sessions/pki.tex}{% Illustrative example (hand-written); answers vary with the agent and model.
\begin{session}{Example}
\you akeyless-watchover pki internal TLS for *.payments.internal, 90 days, auto-renew

\agent \texttt{model/pki.yaml} with a CA key (\texttt{/pki/ca/internal-root}) and an issuer
(\texttt{/pki/issuers/payments-server}, \texttt{allowed\_domains}, TTL 7776000\,s,
\texttt{auto\_renew} with \texttt{destination\_path}); a requester role for the payments workloads;
an expiry forwarder (\texttt{watchover forwarder --preset certificates}); the CLI to issue and
revoke a certificate.
\end{session}
}{}

\section{\texttt{harden}}\label{skill:harden}
\begin{description}[style=nextline,leftmargin=2em,font=\normalfont\bfseries]
\item[Syntax] \texttt{/\allowbreak{}akeyless-\allowbreak{}watchover harden [current settings or concerns]}
\item[Arguments] [current settings or concerns]
\item[Purpose] Account-level hardening.
\item[Delivers] governance.yaml: account\_settings (Terraform), account\_settings\_cli (CLI), custom fields, key policies, event forwarders.
\item[Example] \texttt{/\allowbreak{}akeyless-\allowbreak{}watchover harden}
\end{description}
\IfFileExists{sessions/harden.tex}{% Illustrative example session (hand-written). See Section \ref{sec:how-examples}.
\begin{session}{Example session}
\you /akeyless-watchover harden

\claude Proposals with reasons, split by where they are delivered: \texttt{account\_settings}
(Terraform: password policy, max versions, JWT TTLs, item deletion protection),
\texttt{account\_settings\_cli} (bound IPs, allowed client types, locks -- applied by
\texttt{watchover account-settings}; \emph{dry-run first}), mandatory custom fields such as
\texttt{owner}, key policies, and critical event forwarders.
\end{session}
}{}

\section{\texttt{gateway}}\label{skill:gateway}
\begin{description}[style=nextline,leftmargin=2em,font=\normalfont\bfseries]
\item[Syntax] \texttt{/\allowbreak{}akeyless-\allowbreak{}watchover gateway \textless{}who administers which Gateway,\allowbreak{} SIEM target\textgreater{}}
\item[Arguments] \textless{}who administers which Gateway, SIEM target\textgreater{}
\item[Purpose] Gateway admin plane and audit-log forwarding.
\item[Delivers] gateway.yaml (allowed\_access, log\_forwarding), gateway\_api setup.
\item[Example] \texttt{/\allowbreak{}akeyless-\allowbreak{}watchover gateway platform admins full,\allowbreak{} PKI team targets only,\allowbreak{} Splunk}
\end{description}
\IfFileExists{sessions/gateway.tex}{% Illustrative example session (hand-written). See Section \ref{sec:how-examples}.
\begin{session}{Example session}
\you /akeyless-watchover gateway platform admins full, PKI team targets only, Splunk

\claude \texttt{model/gateway.yaml}: \texttt{allowed\_access} entries (platform admins through the SSO
auth method with their group sub-claim; PKI team limited to the targets permission),
\texttt{log\_forwarding} of type \texttt{splunk} with \texttt{splunk\_token\_var}, and the reminder to set
\texttt{gateway\_api} -- one state per Gateway cluster.
\end{session}
}{}

\section{\texttt{verify}}\label{skill:verify}
\begin{description}[style=nextline,leftmargin=2em,font=\normalfont\bfseries]
\item[Syntax] \texttt{/\allowbreak{}akeyless-\allowbreak{}watchover verify [persona \textbar{}\allowbreak{} path]}
\item[Arguments] [persona \textbar{} path]
\item[Purpose] Prove who can access what.
\item[Delivers] access\_expectations.yaml checks, \texttt{watchover simulate} / \texttt{watchover verify} commands, Impersonate As steps.
\item[Example] \texttt{/\allowbreak{}akeyless-\allowbreak{}watchover verify can the checkout pods read refunds prod?}
\end{description}
\IfFileExists{sessions/verify.tex}{% Illustrative example session (hand-written). See Section \ref{sec:how-examples}.
\begin{session}{Example session}
\you akeyless-watchover verify can the checkout pods read refunds prod?

\agent \textbf{No.} \texttt{checkout-prod-workload} only matches \texttt{/payments/checkout/prod/*}; no
other role is associated with the K8s auth method. Adds the negative check
\texttt{\{ persona: checkout-prod-pod, path: /payments/refunds/prod/db, expect: [] \}} and shows how
to prove it offline (\texttt{watchover simulate --explain}) and live (\texttt{watchover verify}, or
\emph{Impersonate As} in the console).
\end{session}
}{}

\section{\texttt{cli}}\label{skill:cli}
\begin{description}[style=nextline,leftmargin=2em,font=\normalfont\bfseries]
\item[Syntax] \texttt{/\allowbreak{}akeyless-\allowbreak{}watchover cli \textless{}task\textgreater{}}
\item[Arguments] \textless{}task\textgreater{}
\item[Purpose] Exact Akeyless CLI commands for a task (break-glass, bootstrap, one-off checks).
\item[Delivers] Commands checked against the CLI reference, with the equivalent model change.
\item[Example] \texttt{/\allowbreak{}akeyless-\allowbreak{}watchover cli create a temporary read rule for 60 minutes}
\end{description}
\IfFileExists{sessions/cli.tex}{% Illustrative example session (hand-written). See Section \ref{sec:how-examples}.
\begin{session}{Example session}
\you akeyless-watchover cli create a temporary read rule for 60 minutes

\agent \begin{lstlisting}
akeyless set-role-rule --role-name /governance/roles/payments/team-admin \
  --path "/payments/refunds/prod/*" --capability read --capability list --ttl 60
\end{lstlisting}
With the caveat that temporary rules are CLI/API-only and the next apply on that role removes
them -- and the model change to make if the access should last.
\end{session}
}{}

\section{\texttt{explain}}\label{skill:explain}
\begin{description}[style=nextline,leftmargin=2em,font=\normalfont\bfseries]
\item[Syntax] \texttt{/\allowbreak{}akeyless-\allowbreak{}watchover explain \textless{}topic\textgreater{}}
\item[Arguments] \textless{}topic\textgreater{}
\item[Purpose] Explain an Akeyless concept precisely (sub-claims, deny, wildcards, anti-escalation...).
\item[Delivers] A short answer with the verified facts and VERIFY caveats.
\item[Example] \texttt{/\allowbreak{}akeyless-\allowbreak{}watchover explain should this run through the API or the CLI?}
\end{description}
\IfFileExists{sessions/explain.tex}{% Illustrative example session (hand-written). See Section \ref{sec:how-examples}.
\begin{session}{Example session}
\you /akeyless-watchover explain how do sub-claims combine?

\claude Within one key the listed values are alternatives (OR); across keys all must match (AND). With
\texttt{force\_sub\_claims} on the auth method, an association without sub-claims is rejected --
the kit's guard enforces it before Akeyless does. Glob patterns in values are supported; case
sensitivity is per association.
\end{session}
}{}

\chapter{Kit commands (\texttt{scripts/watchover})}\label{ch:kit-commands}
The kit runs in your governance repository: \texttt{scripts/watchover <command> [options]}.
\texttt{scripts/watchover help} prints the same information in the terminal and
\texttt{scripts/watchover <command> --help} shows one command. Each section lists \emph{all} options:
\textbf{required} options must be given, \textbf{repeatable} options may be given several times,
and a default is shown when the option has one. Each section ends with a real execution,
captured when this manual was built (see Section~\ref{sec:how-examples}).

\begin{center}\small
\begin{tabularx}{\linewidth}{@{}l>{\raggedright\arraybackslash}Xl@{}}
\toprule
\textbf{Command} & \textbf{Summary} & \textbf{Akeyless} \\ \midrule
\hyperref[kit:test]{\texttt{test}} & Run every offline check (what CI runs). & offline \\
\hyperref[kit:simulate]{\texttt{simulate}} & Offline: evaluate access\_expectations.yaml against the model. & offline \\
\hyperref[kit:verify]{\texttt{verify}} & Live: compare expectations with describe-permissions run as each persona's CLI profile. & live \\
\hyperref[kit:export]{\texttt{export}} & Export existing config and translate it: YAML model + import blocks, or plain Terraform (-{}-hcl). & live \\
\hyperref[kit:forwarder]{\texttt{forwarder}} & Scaffold a new event forwarder (model entry or HCL); -{}-import-existing adopts one by name. & offline \\
\hyperref[kit:conjur]{\texttt{conjur}} & Translate Conjur policy into model roles + identity bindings. & offline \\
\hyperref[kit:account-settings]{\texttt{account-settings}} & Apply account\_settings\_cli (settings the provider cannot manage) with the CLI. & live \\
\hyperref[kit:token]{\texttt{token}} & CI: exchange the pipeline JWT for an Akeyless token. & live \\
\hyperref[kit:build-skill]{\texttt{build-skill}} & Build the agent-neutral skill package (Agent Skills ZIP) and the prompt pack. & offline \\
\hyperref[kit:install-skill]{\texttt{install-skill}} & Install the skill where an AI agent looks for skills (per-agent folders or any folder). & offline \\
\hyperref[kit:validate-skill]{\texttt{validate-skill}} & Check a skill folder against the Agent Skills specification. & offline \\
\hyperref[kit:mcp]{\texttt{mcp}} & Serve the kit, the docs and the skill commands to any MCP client (stdio). & offline \\
\hyperref[kit:help]{\texttt{help}} & Help in the terminal; \texttt{--markdown} regenerates docs/usage-guide.md. & offline \\
\hyperref[kit:version]{\texttt{version}} & Print the version. & offline \\
\bottomrule
\end{tabularx}
\end{center}

\section{\texttt{watchover test}}\label{kit:test}
Run every offline check (what CI runs). \emph{Offline.}

\begin{Verbatim}[fontsize=\footnotesize,breaklines,frame=leftline,framerule=1pt,rulecolor=\color{wo-rule}]
Run the complete offline test suite (no Akeyless account needed): fmt, validate, terraform test, guard messages, export round trip, bootstrap tests, access simulation, script tests, guide drift check.
\end{Verbatim}

\paragraph{Options}\leavevmode
{\small
\begin{longtable}{@{}>{\raggedright\arraybackslash}p{0.27\linewidth}>{\raggedright\arraybackslash}p{0.17\linewidth}>{\raggedright\arraybackslash}p{0.13\linewidth}>{\raggedright\arraybackslash}p{0.35\linewidth}@{}}
\toprule
\textbf{Option} & \textbf{Value} & \textbf{Required / default} & \textbf{Description} \\ \midrule\endhead
\texttt{-\allowbreak{}{}-\allowbreak{}tool} & \texttt{terraform\textbar{}\allowbreak{}tofu} & default terraform & IaC binary to use \\
\texttt{-\allowbreak{}{}-\allowbreak{}skip-\allowbreak{}iac} & -- & optional & only run the Python checks \\
\bottomrule
\end{longtable}}

\paragraph{Usage patterns}\leavevmode
\begin{Verbatim}[fontsize=\footnotesize,breaklines,breakanywhere]
scripts/watchover test
scripts/watchover test --tool tofu
\end{Verbatim}
\subsection*{Example execution}\addcontentsline{toc}{subsection}{Example execution}
% GENERATED by docs/manual/build.py - do not edit; change the source it names instead.
% source: docs/manual/build.py EXAMPLES (executed in a sandbox copy of the repository)
Every offline check, as CI runs them (formatting, validation, Terraform tests with a mocked provider, guard messages, export round trip, bootstrap tests, Python tests, access simulation):
\begin{terminal}
\begin{Verbatim}[fontsize=\scriptsize,breaklines,breakanywhere,commandchars=\\\{\}]
\textcolor{wo-prompt}{\$} \textbf{scripts/watchover test --tool tofu}

$ tofu fmt -check -recursive   (in terraform)

$ tofu init -backend=false -input=false   (in terraform)

Initializing provider plugins...
- terraform.io/builtin/terraform is built in to OpenTofu
- Reusing previous version of akeyless-community/akeyless from the dependency lock file
- Installing akeyless-community/akeyless v2.0.5...
- Installed akeyless-community/akeyless v2.0.5 (unauthenticated)

OpenTofu has been successfully initialized!

$ tofu validate   (in terraform)
Success! The configuration is valid.

$ tofu test   (in terraform)
tests/model.tftest.hcl... pass
  run "reference_model_plans"... pass
  run "gateway_objects_need_gateway_api"... pass
  run "ara_rule_allowed_when_enabled"... pass
  run "root_scope_allowed_with_flag"... pass
  run "dup"... pass
  run "deny_mixed"... pass
  run "bad_capability"... pass
  run "bad_rule_type"... pass
  run "sra_conflict"... pass
  run "root_scope"... pass
  run "unknown_am"... pass
  run "missing_subclaims"... pass
  run "bad_prefix"... pass
  run "bad_mgmt"... pass
  run "bad_issuer"... pass
  run "bad_gateway"... pass
  run "bad_event"... pass
  run "scaffolded_forwarders_plan"... pass
  run "bad_periodic"... pass
  run "slack_no_url"... pass
  run "empty_gateway_api_is_unset"... pass
  run "naming_check_disabled_with_empty_string"... pass
[... 37 lines omitted ...]
Success! The configuration is valid.

$ tofu test   (in bootstrap)
tests/bootstrap.tftest.hcl... pass
  run "github"... pass
  run "gitlab"... pass
  run "bad_platform"... pass

Success! 3 passed, 0 failed.

$ /usr/local/bin/python3 ./scripts/akl_simulate.py --model terraform/model --expectations terraform/model/access_expectations.yaml   (in .)
PASS  payments-admin           item        /payments/checkout/nonprod/db/app                  expect=['create', 'delete', 'list', 'read', 'update'] got=['create', 'delete', 'list', 'read', 'update']
PASS  payments-admin           item        /payments/checkout/prod/db/app                     expect=['list', 'read'] got=['list', 'read']
PASS  payments-admin           item        /payments/checkout/prod/break-glass/root           expect=[] got=[]
PASS  payments-admin           item        /hr/payroll/prod/db                                expect=[] got=[]
PASS  auditor                  item        /payments/checkout/prod/db/app                     expect=['list'] got=['list']
PASS  auditor                  role        /governance/roles/payments/team-admin              expect=['list', 'read'] got=['list', 'read']
PASS  random-employee          item        /payments/checkout/prod/db/app                     expect=[] got=[]
PASS  checkout-prod-pod        item        /payments/checkout/prod/db/app                     expect=['list', 'read'] got=['list', 'read']
PASS  checkout-prod-pod        item        /payments/checkout/prod/break-glass/root           expect=[] got=[]
PASS  checkout-prod-pod        item        /payments/checkout/nonprod/db/app                  expect=[] got=[]
PASS  checkout-prod-pod        item        /pki/issuers/payments-server                       expect=['list', 'read'] got=['list', 'read']

0 mismatch(es)

$ /usr/local/bin/python3 -m pytest -q scripts/tests   (in .)
.......................................................................                      [100%]
71 passed in <t>s

all checks passed
\end{Verbatim}
\exitcode{0}
\end{terminal}

Only the Python checks:
\begin{terminal}
\begin{Verbatim}[fontsize=\scriptsize,breaklines,breakanywhere,commandchars=\\\{\}]
\textcolor{wo-prompt}{\$} \textbf{scripts/watchover test --skip-iac}

$ /usr/local/bin/python3 ./scripts/akl_simulate.py --model terraform/model --expectations terraform/model/access_expectations.yaml   (in .)
PASS  payments-admin           item        /payments/checkout/nonprod/db/app                  expect=['create', 'delete', 'list', 'read', 'update'] got=['create', 'delete', 'list', 'read', 'update']
PASS  payments-admin           item        /payments/checkout/prod/db/app                     expect=['list', 'read'] got=['list', 'read']
PASS  payments-admin           item        /payments/checkout/prod/break-glass/root           expect=[] got=[]
PASS  payments-admin           item        /hr/payroll/prod/db                                expect=[] got=[]
PASS  auditor                  item        /payments/checkout/prod/db/app                     expect=['list'] got=['list']
PASS  auditor                  role        /governance/roles/payments/team-admin              expect=['list', 'read'] got=['list', 'read']
PASS  random-employee          item        /payments/checkout/prod/db/app                     expect=[] got=[]
PASS  checkout-prod-pod        item        /payments/checkout/prod/db/app                     expect=['list', 'read'] got=['list', 'read']
PASS  checkout-prod-pod        item        /payments/checkout/prod/break-glass/root           expect=[] got=[]
PASS  checkout-prod-pod        item        /payments/checkout/nonprod/db/app                  expect=[] got=[]
PASS  checkout-prod-pod        item        /pki/issuers/payments-server                       expect=['list', 'read'] got=['list', 'read']

0 mismatch(es)

$ /usr/local/bin/python3 -m pytest -q scripts/tests   (in .)
.......................................................................                      [100%]
71 passed in <t>s

all checks passed
\end{Verbatim}
\exitcode{0}
\end{terminal}



\section{\texttt{watchover simulate}}\label{kit:simulate}
Offline: evaluate access\_expectations.yaml against the model. \emph{Offline.}

\begin{Verbatim}[fontsize=\footnotesize,breaklines,frame=leftline,framerule=1pt,rulecolor=\color{wo-rule}]
Offline access check of the YAML model against expected access (no Akeyless calls).

Usage:
  python3 scripts/akl_simulate.py --model terraform/model       --expectations terraform/model/access_expectations.yaml [--explain]

Exit code 0 = all expectations met, 1 = mismatch, 2 = usage/model error.
\end{Verbatim}

\paragraph{Options}\leavevmode
{\small
\begin{longtable}{@{}>{\raggedright\arraybackslash}p{0.27\linewidth}>{\raggedright\arraybackslash}p{0.17\linewidth}>{\raggedright\arraybackslash}p{0.13\linewidth}>{\raggedright\arraybackslash}p{0.35\linewidth}@{}}
\toprule
\textbf{Option} & \textbf{Value} & \textbf{Required / default} & \textbf{Description} \\ \midrule\endhead
\texttt{-\allowbreak{}{}-\allowbreak{}model} & \texttt{MODEL} & \textbf{required} & model directory (terraform/model) \\
\texttt{-\allowbreak{}{}-\allowbreak{}expectations} & \texttt{EXPECTATIONS} & \textbf{required} & access\_expectations.yaml \\
\texttt{-\allowbreak{}{}-\allowbreak{}explain} & -- & optional & print the matching rules for every check, not only failures \\
\bottomrule
\end{longtable}}

\paragraph{Usage patterns}\leavevmode
\begin{Verbatim}[fontsize=\footnotesize,breaklines,breakanywhere]
scripts/watchover simulate --model terraform/model --expectations terraform/model/access_expectations.yaml --explain
\end{Verbatim}
\subsection*{Example execution}\addcontentsline{toc}{subsection}{Example execution}
% GENERATED by docs/manual/build.py - do not edit; change the source it names instead.
% source: docs/manual/build.py EXAMPLES (executed in a sandbox copy of the repository)
Offline evaluation of the access expectations against the model:
\begin{terminal}
\begin{Verbatim}[fontsize=\scriptsize,breaklines,breakanywhere,commandchars=\\\{\}]
\textcolor{wo-prompt}{\$} \textbf{scripts/watchover simulate --model terraform/model --expectations terraform/model/access_expectations.yaml}
PASS  payments-admin           item        /payments/checkout/nonprod/db/app                  expect=['create', 'delete', 'list', 'read', 'update'] got=['create', 'delete', 'list', 'read', 'update']
PASS  payments-admin           item        /payments/checkout/prod/db/app                     expect=['list', 'read'] got=['list', 'read']
PASS  payments-admin           item        /payments/checkout/prod/break-glass/root           expect=[] got=[]
PASS  payments-admin           item        /hr/payroll/prod/db                                expect=[] got=[]
PASS  auditor                  item        /payments/checkout/prod/db/app                     expect=['list'] got=['list']
PASS  auditor                  role        /governance/roles/payments/team-admin              expect=['list', 'read'] got=['list', 'read']
PASS  random-employee          item        /payments/checkout/prod/db/app                     expect=[] got=[]
PASS  checkout-prod-pod        item        /payments/checkout/prod/db/app                     expect=['list', 'read'] got=['list', 'read']
PASS  checkout-prod-pod        item        /payments/checkout/prod/break-glass/root           expect=[] got=[]
PASS  checkout-prod-pod        item        /payments/checkout/nonprod/db/app                  expect=[] got=[]
PASS  checkout-prod-pod        item        /pki/issuers/payments-server                       expect=['list', 'read'] got=['list', 'read']

0 mismatch(es)
\end{Verbatim}
\exitcode{0}
\end{terminal}

With \texttt{-\allowbreak{}{}-\allowbreak{}explain}, every check shows the rules that produced the result:
\begin{terminal}
\begin{Verbatim}[fontsize=\scriptsize,breaklines,breakanywhere,commandchars=\\\{\}]
\textcolor{wo-prompt}{\$} \textbf{scripts/watchover simulate --model terraform/model --expectations terraform/model/access_expectations.yaml --explain | head -n 12}
PASS  payments-admin           item        /payments/checkout/nonprod/db/app                  expect=['create', 'delete', 'list', 'read', 'update'] got=['create', 'delete', 'list', 'read', 'update']
        /governance/roles/payments/team-admin: /payments/+/nonprod/* -> ['create', 'delete', 'list', 'read', 'update']
PASS  payments-admin           item        /payments/checkout/prod/db/app                     expect=['list', 'read'] got=['list', 'read']
        /governance/roles/payments/team-admin: /payments/+/prod/* -> ['list', 'read']
PASS  payments-admin           item        /payments/checkout/prod/break-glass/root           expect=[] got=[]
        /governance/roles/payments/team-admin: /payments/+/prod/* -> ['list', 'read']
        /governance/roles/payments/team-admin: /payments/+/prod/break-glass/* -> ['deny']
        deny overrides all other permissions
PASS  payments-admin           item        /hr/payroll/prod/db                                expect=[] got=[]
        no matching rule
PASS  auditor                  item        /payments/checkout/prod/db/app                     expect=['list'] got=['list']
        /governance/roles/platform/auditors: /* -> ['list']
\end{Verbatim}
\exitcode{0}
\end{terminal}



\section{\texttt{watchover verify}}\label{kit:verify}
Live: compare expectations with describe-permissions run as each persona's CLI profile. \emph{Calls Akeyless.}

\begin{Verbatim}[fontsize=\footnotesize,breaklines,frame=leftline,framerule=1pt,rulecolor=\color{wo-rule}]
Live access verification against Akeyless (after apply).

For every check in the expectations file, runs - as the persona's CLI profile -
    akeyless describe-permissions --path <path> --type <item|am|role|target> --json [--profile P]
which returns the *caller's* own permissions on that path ({"client_permissions": [...]}),
and compares them with the expected set.

Personas need a `profile:` (an `akeyless configure` profile authenticating as that
persona, e.g. a test SSO user in the right IdP group or a test workload identity).
Personas without a profile are skipped (use the Console "Impersonate As" for those).

Usage:
  python3 scripts/akl_verify.py --expectations terraform/model/access_expectations.yaml [--akeyless akeyless]
Exit code 0 = all checked expectations met, 1 = mismatch, 2 = execution error.
\end{Verbatim}

\paragraph{Options}\leavevmode
{\small
\begin{longtable}{@{}>{\raggedright\arraybackslash}p{0.27\linewidth}>{\raggedright\arraybackslash}p{0.17\linewidth}>{\raggedright\arraybackslash}p{0.13\linewidth}>{\raggedright\arraybackslash}p{0.35\linewidth}@{}}
\toprule
\textbf{Option} & \textbf{Value} & \textbf{Required / default} & \textbf{Description} \\ \midrule\endhead
\texttt{-\allowbreak{}{}-\allowbreak{}expectations} & \texttt{EXPECTATIONS} & \textbf{required} & access\_expectations.yaml (personas need a 'profile') \\
\texttt{-\allowbreak{}{}-\allowbreak{}akeyless} & \texttt{AKEYLESS} & default akeyless & path to the akeyless CLI binary \\
\texttt{-\allowbreak{}{}-\allowbreak{}error-\allowbreak{}means-\allowbreak{}no-\allowbreak{}access} & -- & optional & treat a failing describe-permissions call as 'no permissions' (check the CLI behaviour first) \\
\bottomrule
\end{longtable}}

\paragraph{Usage patterns}\leavevmode
\begin{Verbatim}[fontsize=\footnotesize,breaklines,breakanywhere]
scripts/watchover verify --expectations terraform/model/access_expectations.yaml
\end{Verbatim}
\subsection*{Example execution}\addcontentsline{toc}{subsection}{Example execution}
% GENERATED by docs/manual/build.py - do not edit; change the source it names instead.
% source: docs/manual/build.py EXAMPLES (executed in a sandbox copy of the repository)
The stand-in account has one console change that is not in Git:
\begin{terminal}
\begin{Verbatim}[fontsize=\scriptsize,breaklines,breakanywhere,commandchars=\\\{\}]
\textcolor{wo-prompt}{\$} \textbf{scripts/watchover verify --expectations terraform/model/access_expectations.yaml}
PASS  payments-admin           /payments/checkout/nonprod/db/app                  expect=['create', 'delete', 'list', 'read', 'update'] got=['create', 'delete', 'list', 'read', 'update']
PASS  payments-admin           /payments/checkout/prod/db/app                     expect=['list', 'read'] got=['list', 'read']
PASS  payments-admin           /payments/checkout/prod/break-glass/root           expect=[] got=[]
PASS  payments-admin           /hr/payroll/prod/db                                expect=[] got=[]
PASS  auditor                  /payments/checkout/prod/db/app                     expect=['list'] got=['list']
[... 1 lines omitted ...]
FAIL  random-employee          /payments/checkout/prod/db/app                     expect=[] got=['list', 'read']
PASS  checkout-prod-pod        /payments/checkout/prod/db/app                     expect=['list', 'read'] got=['list', 'read']
PASS  checkout-prod-pod        /payments/checkout/prod/break-glass/root           expect=[] got=[]
PASS  checkout-prod-pod        /payments/checkout/nonprod/db/app                  expect=[] got=[]
PASS  checkout-prod-pod        /pki/issuers/payments-server                       expect=['list', 'read'] got=['list', 'read']

1 mismatch(es), 0 check(s) skipped (persona without profile)
\end{Verbatim}
\exitcode{1}
\end{terminal}



\section{\texttt{watchover export}}\label{kit:export}
Export existing config and translate it: YAML model + import blocks, or plain Terraform (-{}-hcl). \emph{Calls Akeyless.}

\begin{Verbatim}[fontsize=\footnotesize,breaklines,frame=leftline,framerule=1pt,rulecolor=\color{wo-rule}]
Export an existing Akeyless account and translate it into Terraform.

Two output modes (both read the same objects with the akeyless CLI, or from saved JSON):

  model (default)  YAML model for terraform/ + import blocks for its addresses
                   <out>/model/roles/imported.yaml          roles, rules, management plane, associations
                   <out>/model/auth_methods.imported.yaml   auth methods incl. type-specific settings
                   <out>/model/folders.imported.yaml        explicit folders (description, tags, protection)
                   <out>/model/governance.imported.yaml     custom fields, key policies
                   <out>/imports.tf                         import {} blocks for terraform/*.tf
  hcl (--hcl)      Plain Terraform generated by Terraform itself from the live objects
                   <out>/hcl/providers.tf + imports.tf      flat import blocks for EVERY importable object found
                                                            (also PKI issuers and account settings)
                   <out>/hcl/generated.tf                   with --generate: `terraform plan -generate-config-out`

Always written: <out>/raw/ (every JSON read, secrets redacted - reusable with --from-dir)
and <out>/REPORT.md (everything that needs a human decision).

Selection: --types limits object families; --only keeps names matching a regex;
--persona-auth-method + --persona-claim keep one identity's view (its auth method and the roles
whose association sub-claims match it - e.g. "the OIDC config and access of one user").

Usage:
  watchover export --out export/ --profile admin
  watchover export --out export/ --profile admin --types auth-methods --only 'oidc'
  watchover export --out export/ --profile admin --persona-auth-method /governance/auth/sso-entra-oidc \
      --persona-claim email=alice@example.com --persona-claim groups=grp-akl-payments-admins
  watchover export --out export/ --profile admin --hcl --generate      # needs terraform + AKEYLESS_ACCESS_ID/KEY
  watchover export --out export/ --from-dir export-old/raw              # offline replay
\end{Verbatim}

\paragraph{Options}\leavevmode
{\small
\begin{longtable}{@{}>{\raggedright\arraybackslash}p{0.27\linewidth}>{\raggedright\arraybackslash}p{0.17\linewidth}>{\raggedright\arraybackslash}p{0.13\linewidth}>{\raggedright\arraybackslash}p{0.35\linewidth}@{}}
\toprule
\textbf{Option} & \textbf{Value} & \textbf{Required / default} & \textbf{Description} \\ \midrule\endhead
\texttt{-\allowbreak{}{}-\allowbreak{}out} & \texttt{OUT} & \textbf{required} & output directory (model/, imports.tf or hcl/, raw/, REPORT.md) \\
\texttt{-\allowbreak{}{}-\allowbreak{}transport} & \texttt{cli\textbar{}\allowbreak{}api} & default cli & how to read Akeyless: the akeyless CLI (profiles, SSO logins) or the REST API directly (token) \\
\texttt{-\allowbreak{}{}-\allowbreak{}profile} & \texttt{PROFILE} & optional & cli transport: akeyless CLI profile with read access to the objects to export \\
\texttt{-\allowbreak{}{}-\allowbreak{}akeyless} & \texttt{AKEYLESS} & default akeyless & cli transport: path to the akeyless CLI binary \\
\texttt{-\allowbreak{}{}-\allowbreak{}gateway-\allowbreak{}url} & \texttt{GATEWAY\_\allowbreak{}URL} & optional & cli transport, -{}-types gateway: Gateway configuration URL (https://gw:8000) \\
\texttt{-\allowbreak{}{}-\allowbreak{}api} & \texttt{API} & default https://api.akeyless.io & api transport: API base URL (default \$AKEYLESS\_API or https://api.akeyless.io) \\
\texttt{-\allowbreak{}{}-\allowbreak{}token} & \texttt{TOKEN} & optional & api transport: Akeyless token (default \$AKEYLESS\_TOKEN; e.g. from 'watchover token') \\
\texttt{-\allowbreak{}{}-\allowbreak{}gateway-\allowbreak{}api} & \texttt{GATEWAY\_\allowbreak{}API} & optional & api transport, -{}-types gateway: Gateway V2 API (https://gw:8000/api/v2) \\
\texttt{-\allowbreak{}{}-\allowbreak{}from-\allowbreak{}dir} & \texttt{FROM\_\allowbreak{}DIR} & optional & replay a raw/ directory of a previous export instead of calling the CLI \\
\texttt{-\allowbreak{}{}-\allowbreak{}types} & \texttt{TYPES} & default roles,auth-methods,folders,custom-fields,key-policies,pki-issuers,account-settings,targets,groups & comma-separated object families (default: all but gateway): roles, auth-methods, folders, custom-fields, key-policies, pki-issuers, account-settings, targets, groups, gateway \\
\texttt{-\allowbreak{}{}-\allowbreak{}only} & \texttt{ONLY} & optional & keep only objects whose name matches this regex \\
\texttt{-\allowbreak{}{}-\allowbreak{}exclude} & \texttt{EXCLUDE} & \emph{repeatable} & regex of names to skip (repeatable) \\
\texttt{-\allowbreak{}{}-\allowbreak{}root} & \texttt{ROOT} & default / & folder to start the item/folder walk from \\
\texttt{-\allowbreak{}{}-\allowbreak{}max-\allowbreak{}depth} & \texttt{MAX\_\allowbreak{}DEPTH} & default 20 & maximum folder depth of the walk \\
\texttt{-\allowbreak{}{}-\allowbreak{}persona-\allowbreak{}auth-\allowbreak{}method} & \texttt{PERSONA\_\allowbreak{}AUTH\_\allowbreak{}METHOD} & optional & keep one identity's view: its auth method ... \\
\texttt{-\allowbreak{}{}-\allowbreak{}persona-\allowbreak{}claim} & \texttt{PERSONA\_\allowbreak{}CLAIM} & \emph{repeatable} & ... and its claims, key=value (repeatable); keeps roles whose association matches \\
\texttt{-\allowbreak{}{}-\allowbreak{}hcl} & -- & optional & write plain Terraform import blocks for every object (hcl/) \\
\texttt{-\allowbreak{}{}-\allowbreak{}generate} & -- & optional & with -{}-hcl: run 'terraform plan -generate-config-out' (needs terraform and AKEYLESS\_ACCESS\_ID/AKEYLESS\_ACCESS\_KEY in the environment) \\
\texttt{-\allowbreak{}{}-\allowbreak{}terraform} & \texttt{TERRAFORM} & default terraform & terraform or tofu binary used by -{}-generate \\
\bottomrule
\end{longtable}}

\paragraph{Usage patterns}\leavevmode
\begin{Verbatim}[fontsize=\footnotesize,breaklines,breakanywhere]
scripts/watchover export --out export/ --profile admin --exclude '^/governance/auth/ci-terraform$'
scripts/watchover export --out export/ --profile admin --types auth-methods --only sso
scripts/watchover export --out export/ --profile admin --persona-auth-method /governance/auth/sso-entra-oidc --persona-claim email=alice@example.com --persona-claim groups=grp-akl-payments-admins
AKEYLESS_ACCESS_ID=p-... AKEYLESS_ACCESS_KEY=... scripts/watchover export --out export/ --profile admin --hcl --generate
AKEYLESS_TOKEN=$(scripts/watchover token) scripts/watchover export --out export/ --transport api --hcl --types targets,groups,gateway --gateway-api https://gw.example.com:8000/api/v2
scripts/watchover export --out export2/ --from-dir export/raw     # offline replay
\end{Verbatim}
\subsection*{Example execution}\addcontentsline{toc}{subsection}{Example execution}
% GENERATED by docs/manual/build.py - do not edit; change the source it names instead.
% source: docs/manual/build.py EXAMPLES (executed in a sandbox copy of the repository)
\begin{terminal}
\begin{Verbatim}[fontsize=\scriptsize,breaklines,breakanywhere,commandchars=\\\{\}]
\textcolor{wo-prompt}{\$} \textbf{scripts/watchover export --out export/ --profile admin --types roles,auth-methods && head -n 8 export/imports.tf}
exported (model): 3 roles, 7 auth-methods; 7 item(s) in REPORT.md
# Generated by watchover export (model mode) - addresses of terraform/*.tf.
# Remove after the first successful apply.

import \{
  to = akeyless_auth_method_aws_iam.this["/governance/auth/aws-payments-prod"]
  id = "/governance/auth/aws-payments-prod"
\}
\end{Verbatim}
\exitcode{0}
\end{terminal}



\section{\texttt{watchover forwarder}}\label{kit:forwarder}
Scaffold a new event forwarder (model entry or HCL); -{}-import-existing adopts one by name. \emph{Offline.}

\begin{Verbatim}[fontsize=\footnotesize,breaklines,frame=leftline,framerule=1pt,rulecolor=\color{wo-rule}]
Scaffold Terraform configuration for Akeyless event forwarders.

Akeyless has no operation that lists event forwarders, so they can't be exported. This command
writes new ones instead, already validated against the same rules the Terraform guards enforce:

  * model entry for terraform/model/governance.yaml (default; --write inserts it in place)
  * plain HCL resource block (--format hcl) for repositories that don't use the YAML model
  * an import block (--import-existing) to bring an EXISTING forwarder whose name you know under
    Terraform: the provider imports event forwarders by name

Secrets (webhook passwords/tokens/certificates, Slack/Teams URLs, ServiceNow credentials) are
never written: the output references a *_var key you then provide in TF_VAR_event_forwarder_secrets.

Examples:
  watchover forwarder --name /governance/events/secops-mail --type email --email-to secops@example.com \
      --preset critical --items-source '/*'
  watchover forwarder --name /governance/events/pki-slack --type slack --preset certificates \
      --items-source '/pki/*' --write terraform/model/governance.yaml
  watchover forwarder --name /governance/events/siem --type webhook --url https://siem.example.com/akeyless \
      --auth-type bearer-token --preset critical --auth-method-source '/*' --format hcl
  watchover forwarder --name /governance/events/old-mail --type email --email-to a@b.c --event certificate-expired \
      --import-existing
\end{Verbatim}

\paragraph{Options}\leavevmode
{\small
\begin{longtable}{@{}>{\raggedright\arraybackslash}p{0.27\linewidth}>{\raggedright\arraybackslash}p{0.17\linewidth}>{\raggedright\arraybackslash}p{0.13\linewidth}>{\raggedright\arraybackslash}p{0.35\linewidth}@{}}
\toprule
\textbf{Option} & \textbf{Value} & \textbf{Required / default} & \textbf{Description} \\ \midrule\endhead
\texttt{-\allowbreak{}{}-\allowbreak{}name} & \texttt{NAME} & \textbf{required} & forwarder name, e.g. /governance/events/secops-mail \\
\texttt{-\allowbreak{}{}-\allowbreak{}type} & \texttt{email\textbar{}\allowbreak{}webhook\textbar{}\allowbreak{}slack\textbar{}\allowbreak{}teams\textbar{}\allowbreak{}service\_\allowbreak{}now} & \textbf{required} & forwarder type \\
\texttt{-\allowbreak{}{}-\allowbreak{}description} & \texttt{DESCRIPTION} & optional & description \\
\texttt{-\allowbreak{}{}-\allowbreak{}preset} & \texttt{access-\allowbreak{}requests\textbar{}\allowbreak{}auth-\allowbreak{}methods\textbar{}\allowbreak{}certificates\textbar{}\allowbreak{}critical\textbar{}\allowbreak{}operations\textbar{}\allowbreak{}rotation\textbar{}\allowbreak{}secrets} & \emph{repeatable} & add a group of event types (repeatable); see 'Presets' in the usage guide \\
\texttt{-\allowbreak{}{}-\allowbreak{}event} & \texttt{EVENT} & \emph{repeatable} & add one event type (repeatable) \\
\texttt{-\allowbreak{}{}-\allowbreak{}items-\allowbreak{}source} & \texttt{ITEMS\_\allowbreak{}SOURCE} & \emph{repeatable} & item path pattern to watch, e.g. '/pki/*' (repeatable) \\
\texttt{-\allowbreak{}{}-\allowbreak{}targets-\allowbreak{}source} & \texttt{TARGETS\_\allowbreak{}SOURCE} & \emph{repeatable} & target path pattern to watch (repeatable) \\
\texttt{-\allowbreak{}{}-\allowbreak{}auth-\allowbreak{}method-\allowbreak{}source} & \texttt{AUTH\_\allowbreak{}METHOD\_\allowbreak{}SOURCE} & \emph{repeatable} & auth method path pattern to watch (repeatable) \\
\texttt{-\allowbreak{}{}-\allowbreak{}gateway-\allowbreak{}source} & \texttt{GATEWAY\_\allowbreak{}SOURCE} & \emph{repeatable} & Gateway cluster URL to watch (repeatable; required for teams) \\
\texttt{-\allowbreak{}{}-\allowbreak{}runner} & \texttt{immediate\textbar{}\allowbreak{}periodic} & default immediate & send at once or as a periodic digest \\
\texttt{-\allowbreak{}{}-\allowbreak{}every} & \texttt{EVERY} & optional & hours between digests (required with -{}-runner periodic) \\
\texttt{-\allowbreak{}{}-\allowbreak{}disabled} & -- & optional & create it disabled \\
\texttt{-\allowbreak{}{}-\allowbreak{}key} & \texttt{KEY} & optional & protection key for the forwarder's secret (default: account key) \\
\texttt{-\allowbreak{}{}-\allowbreak{}email-\allowbreak{}to} & \texttt{EMAIL\_\allowbreak{}TO} & \emph{repeatable} & email: recipient (repeatable) \\
\texttt{-\allowbreak{}{}-\allowbreak{}include-\allowbreak{}error} & -- & optional & email: include error details \\
\texttt{-\allowbreak{}{}-\allowbreak{}override-\allowbreak{}url} & \texttt{OVERRIDE\_\allowbreak{}URL} & optional & email: Gateway URL used in links (port 18888) \\
\texttt{-\allowbreak{}{}-\allowbreak{}url} & \texttt{URL} & optional & webhook: URL (slack/teams URLs are secrets and go in the secrets map) \\
\texttt{-\allowbreak{}{}-\allowbreak{}auth-\allowbreak{}type} & \texttt{AUTH\_\allowbreak{}TYPE} & optional & webhook: user-pass/bearer-token/certificate; service\_now: user-pass/jwt \\
\texttt{-\allowbreak{}{}-\allowbreak{}username} & \texttt{USERNAME} & optional & webhook: user for user-pass \\
\texttt{-\allowbreak{}{}-\allowbreak{}host} & \texttt{HOST} & optional & service\_now: instance host \\
\texttt{-\allowbreak{}{}-\allowbreak{}admin-\allowbreak{}name} & \texttt{ADMIN\_\allowbreak{}NAME} & optional & service\_now: admin name (user-pass) \\
\texttt{-\allowbreak{}{}-\allowbreak{}client-\allowbreak{}id} & \texttt{CLIENT\_\allowbreak{}ID} & optional & service\_now: client ID (jwt) \\
\texttt{-\allowbreak{}{}-\allowbreak{}user-\allowbreak{}email} & \texttt{USER\_\allowbreak{}EMAIL} & optional & service\_now: user email (jwt) \\
\texttt{-\allowbreak{}{}-\allowbreak{}format} & \texttt{yaml\textbar{}\allowbreak{}hcl} & default yaml & model entry (yaml) or plain resource (hcl) \\
\texttt{-\allowbreak{}{}-\allowbreak{}write} & \texttt{GOVERNANCE\_\allowbreak{}YAML} & optional & yaml: insert the entry into this governance.yaml \\
\texttt{-\allowbreak{}{}-\allowbreak{}import-\allowbreak{}existing} & -- & optional & the forwarder already exists in Akeyless: also print the import block (import ID = name) \\
\bottomrule
\end{longtable}}

\paragraph{Usage patterns}\leavevmode
\begin{Verbatim}[fontsize=\footnotesize,breaklines,breakanywhere]
scripts/watchover forwarder --name /governance/events/secops-mail --type email --email-to secops@example.com --preset critical --items-source '/*' --auth-method-source '/*'
scripts/watchover forwarder --name /governance/events/pki-slack --type slack --preset certificates --items-source '/pki/*' --write terraform/model/governance.yaml
scripts/watchover forwarder --name /governance/events/siem --type webhook --url https://siem.example.com/akeyless --auth-type bearer-token --preset critical --auth-method-source '/*' --format hcl
scripts/watchover forwarder --name /governance/events/legacy-mail --type email --email-to ops@example.com --event certificate-expired --items-source '/*' --import-existing
\end{Verbatim}
\subsection*{Example execution}\addcontentsline{toc}{subsection}{Example execution}
% GENERATED by docs/manual/build.py - do not edit; change the source it names instead.
% source: docs/manual/build.py EXAMPLES (executed in a sandbox copy of the repository)
A model entry (default output):
\begin{terminal}
\begin{Verbatim}[fontsize=\scriptsize,breaklines,breakanywhere,commandchars=\\\{\}]
\textcolor{wo-prompt}{\$} \textbf{scripts/watchover forwarder --name /governance/events/secops-mail --type email --email-to secops@example.com --preset critical --items-source '/*' --auth-method-source '/*'}
# add under event_forwarders: in terraform/model/governance.yaml
  /governance/events/secops-mail:
    type: email
    event_types:
    - certificate-expired
    - certificate-provisioning-failure
    - auth-method-expired
    - rotated-secret-failure
    - dynamic-secret-failure
    - multi-auth-failure
    - uid-rotation-failure
    - gateway-inactive
    runner_type: immediate
    items_event_source_locations:
    - /*
    auth_methods_event_source_locations:
    - /*
    email_to:
    - secops@example.com
\end{Verbatim}
\exitcode{0}
\end{terminal}

Insert it into the model (comments and the rest of the file are kept; the secret key is printed): (The sandbox copy of governance.yaml is restored afterwards, so later examples see the original.)
\begin{terminal}
\begin{Verbatim}[fontsize=\scriptsize,breaklines,breakanywhere,commandchars=\\\{\}]
\textcolor{wo-prompt}{\$} \textbf{scripts/watchover forwarder --name /governance/events/pki-slack --type slack --preset certificates --items-source '/pki/*' --runner periodic --every 24 --write terraform/model/governance.yaml && grep -n -A3 'pki-slack' terraform/model/governance.yaml}
added /governance/events/pki-slack to terraform/model/governance.yaml
59:  /governance/events/pki-slack:
60-    type: slack
61-    event_types:
62-    - certificate-pending-expiration

# provide these keys in TF_VAR_event_forwarder_secrets (CI secret, JSON map):
#   governance_events_pki_slack_url
\end{Verbatim}
\exitcode{0}
\end{terminal}

Plain HCL plus an import block for a forwarder that already exists:
\begin{terminal}
\begin{Verbatim}[fontsize=\scriptsize,breaklines,breakanywhere,commandchars=\\\{\}]
\textcolor{wo-prompt}{\$} \textbf{scripts/watchover forwarder --name /governance/events/siem --type webhook --url https://siem.example.com/akeyless --auth-type bearer-token --preset critical --auth-method-source '/*' --format hcl --import-existing}
resource "akeyless_event_forwarder_webhook" "governance_events_siem" \{
  name                                = "/governance/events/siem"
  event_types                         = ["certificate-expired", "certificate-provisioning-failure", "auth-method-expired", "rotated-secret-failure", "dynamic-secret-failure", "multi-auth-failure", "uid-rotation-failure", "gateway-inactive"]
  runner_type                         = "immediate"
  auth_methods_event_source_locations = ["/*"]
  url                                 = "https://siem.example.com/akeyless"
  auth_type                           = "bearer-token"
  auth_token                          = var.event_forwarder_secrets["governance_events_siem_token"]
\}

# imports.tf - the forwarder exists already; Terraform adopts it on the next apply
import \{
  to = akeyless_event_forwarder_webhook.governance_events_siem
  id = "/governance/events/siem"
\}

# needs: variable "event_forwarder_secrets" \{ type = map(string), sensitive = true \} (already declared in terraform/main.tf)

# provide these keys in TF_VAR_event_forwarder_secrets (CI secret, JSON map):
#   governance_events_siem_token
\end{Verbatim}
\exitcode{0}
\end{terminal}

ServiceNow with JWT authentication, created disabled:
\begin{terminal}
\begin{Verbatim}[fontsize=\scriptsize,breaklines,breakanywhere,commandchars=\\\{\}]
\textcolor{wo-prompt}{\$} \textbf{scripts/watchover forwarder --name /governance/events/snow --type service_now --host example.service-now.com --auth-type jwt --client-id snow-client --user-email akeyless@example.com --preset rotation --items-source '/*' --disabled}
# add under event_forwarders: in terraform/model/governance.yaml
  /governance/events/snow:
    type: service_now
    event_types:
    - next-automatic-rotation
    - rotated-secret-success
    - rotated-secret-failure
    - rotation-usage
    runner_type: immediate
    enable: false
    items_event_source_locations:
    - /*
    host: example.service-now.com
    auth_type: jwt
    client_id: snow-client
    user_email: akeyless@example.com
    client_secret_var: governance_events_snow_client_secret
    app_private_key_var: governance_events_snow_app_private_key

# provide these keys in TF_VAR_event_forwarder_secrets (CI secret, JSON map):
#   governance_events_snow_client_secret
#   governance_events_snow_app_private_key
\end{Verbatim}
\exitcode{0}
\end{terminal}

Invalid input is rejected before Terraform runs (exit code 2):
\begin{terminal}
\begin{Verbatim}[fontsize=\scriptsize,breaklines,breakanywhere,commandchars=\\\{\}]
\textcolor{wo-prompt}{\$} \textbf{scripts/watchover forwarder --name /governance/events/ops --type teams --preset operations --items-source '/*'}
error: teams needs --gateway-source (provider 2.0.5 requires gateway_event_source_locations)
\end{Verbatim}
\exitcode{2}
\end{terminal}



\section{\texttt{watchover conjur}}\label{kit:conjur}
Translate Conjur policy into model roles + identity bindings. \emph{Offline.}

\begin{Verbatim}[fontsize=\footnotesize,breaklines,frame=leftline,framerule=1pt,rulecolor=\color{wo-rule}]
Translate CyberArk Conjur policy (RBAC) into the akeyless-watchover YAML model.

Akeyless "Automatic Migration" (Gateway, type conjur) copies *secret values* only.
Conjur's authorization model (policies, hosts, layers, groups, permits, grants) has
to be rebuilt as Akeyless roles - this script does the mechanical part:

  Conjur                               Akeyless model
  -----------------------------------  ------------------------------------------------
  !variable  <policy>/<id>             item  <prefix>/<policy>/<id>
  !permit  privileges on a variable    rule on that item path:
      execute (fetch value)  -> read       read   (+ list implied)
      read    (see metadata) -> list       list
      update  (set value)    -> update     update
  role with permits (layer/group/      one Akeyless role  <role_prefix>conjur/<kind>-<id>
      host/user/policy)
  !grant  role: A, member: B           no role nesting in Akeyless -> B is bound to A's
                                       Akeyless role too (flattened, transitive)
  host / user identities               bindings report: which Akeyless roles each Conjur
                                       identity needs; YOU map it to an auth method +
                                       sub-claims (e.g. authn-jwt -> OAuth2/JWT auth method)

Supported policy statements: !policy (id, body), !variable, !host, !layer, !group, !user,
!webservice, !permit, !grant, !deny (reported, not translated). Annotations are ignored.

Usage:
  python3 scripts/conjur2akeyless.py policy-root.yml [more.yml ...]       --item-prefix /conjur --role-prefix /governance/roles/ --out export/conjur
Writes <out>/roles/conjur.yaml, <out>/bindings.yaml and <out>/REPORT.md.

VERIFY: where the Gateway migration places each Conjur variable (target-location + Conjur
variable id is assumed here) - compare with the migrated item names before applying.
\end{Verbatim}

\paragraph{Options}\leavevmode
{\small
\begin{longtable}{@{}>{\raggedright\arraybackslash}p{0.27\linewidth}>{\raggedright\arraybackslash}p{0.17\linewidth}>{\raggedright\arraybackslash}p{0.13\linewidth}>{\raggedright\arraybackslash}p{0.35\linewidth}@{}}
\toprule
\textbf{Option} & \textbf{Value} & \textbf{Required / default} & \textbf{Description} \\ \midrule\endhead
\texttt{policy} & \texttt{POLICY...} & \textbf{required} & Conjur policy YAML files (loaded at the root namespace in the given order) \\
\texttt{-\allowbreak{}{}-\allowbreak{}item-\allowbreak{}prefix} & \texttt{ITEM\_\allowbreak{}PREFIX} & default /conjur & Akeyless folder the Gateway migration imports into \\
\texttt{-\allowbreak{}{}-\allowbreak{}role-\allowbreak{}prefix} & \texttt{ROLE\_\allowbreak{}PREFIX} & default /governance/roles/ & prefix of the generated Akeyless role names \\
\texttt{-\allowbreak{}{}-\allowbreak{}out} & \texttt{OUT} & \textbf{required} & output directory (roles/conjur.yaml, bindings.yaml, REPORT.md) \\
\bottomrule
\end{longtable}}

\paragraph{Usage patterns}\leavevmode
\begin{Verbatim}[fontsize=\footnotesize,breaklines,breakanywhere]
scripts/watchover conjur root.yml prod.yml --item-prefix /conjur --out export/conjur
\end{Verbatim}
\subsection*{Example execution}\addcontentsline{toc}{subsection}{Example execution}
% GENERATED by docs/manual/build.py - do not edit; change the source it names instead.
% source: docs/manual/build.py EXAMPLES (executed in a sandbox copy of the repository)
\begin{terminal}
\begin{Verbatim}[fontsize=\scriptsize,breaklines,breakanywhere,commandchars=\\\{\}]
\textcolor{wo-prompt}{\$} \textbf{scripts/watchover conjur scripts/tests/fixtures/conjur.yml --out export/conjur && sed -n '1,9p' export/conjur/roles/conjur.yaml}
2 roles, 2 identities, 1 note(s) -> export/conjur
roles:
  /governance/roles/conjur/group-payments-admins:
    allow_external_auth_methods: true
    description: Translated from Conjur group 'payments-admins'
    rules:
    - capabilities:
      - list
      - update
      path: /conjur/prod/payments/db/password
\end{Verbatim}
\exitcode{0}
\end{terminal}



\section{\texttt{watchover account-settings}}\label{kit:account-settings}
Apply account\_settings\_cli (settings the provider cannot manage) with the CLI. \emph{Calls Akeyless.}

\begin{Verbatim}[fontsize=\footnotesize,breaklines,frame=leftline,framerule=1pt,rulecolor=\color{wo-rule}]
Apply the account settings that Terraform provider 2.0.5 cannot manage.

Reads `account_settings_cli` from model/governance.yaml. Keys are the flag names of
`akeyless update-account-settings` (without the leading --):
  * lists  -> the flag is repeated once per value   (--bound-ips a --bound-ips b)
  * bools  -> --flag=true / --flag=false
  * others -> --flag=value
`update-account-settings` is admin-only and idempotent (same input -> same state).

Usage:
  python3 scripts/akl_account_settings.py --governance terraform/model/governance.yaml [--dry-run] [--profile P]
\end{Verbatim}

\paragraph{Options}\leavevmode
{\small
\begin{longtable}{@{}>{\raggedright\arraybackslash}p{0.27\linewidth}>{\raggedright\arraybackslash}p{0.17\linewidth}>{\raggedright\arraybackslash}p{0.13\linewidth}>{\raggedright\arraybackslash}p{0.35\linewidth}@{}}
\toprule
\textbf{Option} & \textbf{Value} & \textbf{Required / default} & \textbf{Description} \\ \midrule\endhead
\texttt{-\allowbreak{}{}-\allowbreak{}governance} & \texttt{GOVERNANCE} & \textbf{required} & path to model/governance.yaml (reads account\_settings\_cli) \\
\texttt{-\allowbreak{}{}-\allowbreak{}profile} & \texttt{PROFILE} & optional & akeyless CLI profile of an account admin \\
\texttt{-\allowbreak{}{}-\allowbreak{}token} & \texttt{TOKEN} & optional & Akeyless token (default: \$AKEYLESS\_TOKEN, e.g. from 'watchover token') \\
\texttt{-\allowbreak{}{}-\allowbreak{}akeyless} & \texttt{AKEYLESS} & default akeyless & path to the akeyless CLI binary \\
\texttt{-\allowbreak{}{}-\allowbreak{}dry-\allowbreak{}run} & -- & optional & print the update-account-settings command, do not run it \\
\bottomrule
\end{longtable}}

\paragraph{Environment variables}\leavevmode\\
\begin{tabularx}{\linewidth}{@{}l>{\raggedright\arraybackslash}X@{}}
\toprule
\texttt{AKEYLESS\_\allowbreak{}TOKEN} & token used when -{}-token is not given \\
\bottomrule
\end{tabularx}

\paragraph{Usage patterns}\leavevmode
\begin{Verbatim}[fontsize=\footnotesize,breaklines,breakanywhere]
scripts/watchover account-settings --governance terraform/model/governance.yaml --dry-run
AKEYLESS_TOKEN=$(scripts/watchover token) scripts/watchover account-settings --governance terraform/model/governance.yaml
\end{Verbatim}
\subsection*{Example execution}\addcontentsline{toc}{subsection}{Example execution}
% GENERATED by docs/manual/build.py - do not edit; change the source it names instead.
% source: docs/manual/build.py EXAMPLES (executed in a sandbox copy of the repository)
\begin{terminal}
\begin{Verbatim}[fontsize=\scriptsize,breaklines,breakanywhere,commandchars=\\\{\}]
\textcolor{wo-prompt}{\$} \textbf{scripts/watchover account-settings --governance terraform/model/governance.yaml --dry-run}
+ akeyless update-account-settings --allowed-client-type=cli --allowed-client-type=ui --allowed-client-type=sdk --bound-ips=10.0.0.0/8 --lock-allowed-client-type=true --lock-bound-ips=true --max-rotation-interval=90 --max-rotation-interval-enable=true --usage-event-enable=true --usage-event-interval=90 --usage-event-object-type=item
\end{Verbatim}
\exitcode{0}
\end{terminal}



\section{\texttt{watchover token}}\label{kit:token}
CI: exchange the pipeline JWT for an Akeyless token. \emph{Calls Akeyless.}

\begin{Verbatim}[fontsize=\footnotesize,breaklines,frame=leftline,framerule=1pt,rulecolor=\color{wo-rule}]
Exchange a CI OIDC/JWT for a short-lived Akeyless token (REST POST /auth, access-type jwt). Prints the token.
\end{Verbatim}

\paragraph{Options}\leavevmode
{\small
\begin{longtable}{@{}>{\raggedright\arraybackslash}p{0.27\linewidth}>{\raggedright\arraybackslash}p{0.17\linewidth}>{\raggedright\arraybackslash}p{0.13\linewidth}>{\raggedright\arraybackslash}p{0.35\linewidth}@{}}
\toprule
\textbf{Option} & \textbf{Value} & \textbf{Required / default} & \textbf{Description} \\ \midrule\endhead
\texttt{-\allowbreak{}{}-\allowbreak{}api} & \texttt{API} & optional & Akeyless API (default: \$AKEYLESS\_API or https://api.akeyless.io) \\
\texttt{-\allowbreak{}{}-\allowbreak{}access-\allowbreak{}id} & \texttt{ACCESS\_\allowbreak{}ID} & optional & access ID of the CI auth method (default: \$AKEYLESS\_ACCESS\_ID) \\
\texttt{-\allowbreak{}{}-\allowbreak{}jwt} & \texttt{JWT} & optional & the CI-issued JWT (default: \$CI\_JWT; prefer the env var - arguments are visible in the process list) \\
\bottomrule
\end{longtable}}

\paragraph{Environment variables}\leavevmode\\
\begin{tabularx}{\linewidth}{@{}l>{\raggedright\arraybackslash}X@{}}
\toprule
\texttt{AKEYLESS\_\allowbreak{}API} & API base URL \\
\texttt{AKEYLESS\_\allowbreak{}ACCESS\_\allowbreak{}ID} & CI auth method access ID \\
\texttt{CI\_\allowbreak{}JWT} & CI-issued JWT \\
\bottomrule
\end{tabularx}

\paragraph{Usage patterns}\leavevmode
\begin{Verbatim}[fontsize=\footnotesize,breaklines,breakanywhere]
CI_JWT=$JWT AKEYLESS_ACCESS_ID=p-xxxx scripts/watchover token
\end{Verbatim}
\subsection*{Example execution}\addcontentsline{toc}{subsection}{Example execution}
% GENERATED by docs/manual/build.py - do not edit; change the source it names instead.
% source: docs/manual/build.py EXAMPLES (executed in a sandbox copy of the repository)
(Shown as you would run it; executed against the local stand-in.)
\begin{terminal}
\begin{Verbatim}[fontsize=\scriptsize,breaklines,breakanywhere,commandchars=\\\{\}]
\textcolor{wo-prompt}{\$} \textbf{CI_JWT=$CI_JWT AKEYLESS_ACCESS_ID=p-ci0demo scripts/watchover token}
t-demo-0b3f9c
\end{Verbatim}
\exitcode{0}
\end{terminal}



\section{\texttt{watchover build-skill}}\label{kit:build-skill}
Build the agent-neutral skill package (Agent Skills ZIP) and the prompt pack. \emph{Offline.}

\begin{Verbatim}[fontsize=\footnotesize,breaklines,frame=leftline,framerule=1pt,rulecolor=\color{wo-rule}]
Build the Agent Skills package (dist/akeyless-watchover/ and .zip) and the prompt pack for assistants without skill support (dist/prompt-pack/ and .zip).
\end{Verbatim}

\paragraph{Options}\leavevmode
{\small
\begin{longtable}{@{}>{\raggedright\arraybackslash}p{0.27\linewidth}>{\raggedright\arraybackslash}p{0.17\linewidth}>{\raggedright\arraybackslash}p{0.13\linewidth}>{\raggedright\arraybackslash}p{0.35\linewidth}@{}}
\toprule
\textbf{Option} & \textbf{Value} & \textbf{Required / default} & \textbf{Description} \\ \midrule\endhead
\texttt{-\allowbreak{}{}-\allowbreak{}out} & \texttt{OUT} & optional & output directory, replaced (default: dist/ in the kit root) \\
\bottomrule
\end{longtable}}

\paragraph{Usage patterns}\leavevmode
\begin{Verbatim}[fontsize=\footnotesize,breaklines,breakanywhere]
scripts/watchover build-skill
scripts/watchover build-skill --out /tmp/wo-dist
\end{Verbatim}
\subsection*{Example execution}\addcontentsline{toc}{subsection}{Example execution}
% GENERATED by docs/manual/build.py - do not edit; change the source it names instead.
% source: docs/manual/build.py EXAMPLES (executed in a sandbox copy of the repository)
Package the skill for upload:
\begin{terminal}
\begin{Verbatim}[fontsize=\scriptsize,breaklines,breakanywhere,commandchars=\\\{\}]
\textcolor{wo-prompt}{\$} \textbf{scripts/watchover build-skill && unzip -l dist/akeyless-watchover.zip | tail -n 1}
built dist/akeyless-watchover.zip (156K)
   337065                     167 files
\end{Verbatim}
\exitcode{0}
\end{terminal}



\section{\texttt{watchover install-skill}}\label{kit:install-skill}
Install the skill where an AI agent looks for skills (per-agent folders or any folder). \emph{Offline.}

\begin{Verbatim}[fontsize=\footnotesize,breaklines,frame=leftline,framerule=1pt,rulecolor=\color{wo-rule}]
Install the skill into the folder an agent scans for skills.
\end{Verbatim}

\paragraph{Options}\leavevmode
{\small
\begin{longtable}{@{}>{\raggedright\arraybackslash}p{0.27\linewidth}>{\raggedright\arraybackslash}p{0.17\linewidth}>{\raggedright\arraybackslash}p{0.13\linewidth}>{\raggedright\arraybackslash}p{0.35\linewidth}@{}}
\toprule
\textbf{Option} & \textbf{Value} & \textbf{Required / default} & \textbf{Description} \\ \midrule\endhead
\texttt{-\allowbreak{}{}-\allowbreak{}agent} & \texttt{agents\textbar{}\allowbreak{}claude\textbar{}\allowbreak{}codex\textbar{}\allowbreak{}copilot\textbar{}\allowbreak{}cursor\textbar{}\allowbreak{}gemini\textbar{}\allowbreak{}custom} & \textbf{required} & target agent; 'agents' = the shared .agents/skills folder; 'custom' needs -{}-dest \\
\texttt{-\allowbreak{}{}-\allowbreak{}scope} & \texttt{project\textbar{}\allowbreak{}user} & default project & project = inside -{}-project; user = your home directory \\
\texttt{-\allowbreak{}{}-\allowbreak{}project} & \texttt{PROJECT} & default . & repository the project-scope skill is installed into \\
\texttt{-\allowbreak{}{}-\allowbreak{}dest} & \texttt{DEST} & optional & explicit skills folder (overrides -{}-agent/-{}-scope) \\
\texttt{-\allowbreak{}{}-\allowbreak{}force} & -- & optional & replace an existing installation \\
\texttt{-\allowbreak{}{}-\allowbreak{}dry-\allowbreak{}run} & -- & optional & print the target folder, change nothing \\
\bottomrule
\end{longtable}}

\paragraph{Usage patterns}\leavevmode
\begin{Verbatim}[fontsize=\footnotesize,breaklines,breakanywhere]
scripts/watchover install-skill --agent codex --project ../akeyless-governance
scripts/watchover install-skill --agent claude --scope user
scripts/watchover install-skill --agent copilot --project . --dry-run
scripts/watchover install-skill --agent custom --dest ~/my-agent/skills --force
\end{Verbatim}
\subsection*{Example execution}\addcontentsline{toc}{subsection}{Example execution}
% GENERATED by docs/manual/build.py - do not edit; change the source it names instead.
% source: docs/manual/build.py EXAMPLES (executed in a sandbox copy of the repository)
\begin{terminal}
\begin{Verbatim}[fontsize=\scriptsize,breaklines,breakanywhere,commandchars=\\\{\}]
\textcolor{wo-prompt}{\$} \textbf{mkdir -p ../governance-repo && scripts/watchover install-skill --agent codex --project ../governance-repo}
installed akeyless-watchover for OpenAI Codex (CLI, IDE, app) in <tmp>/governance-repo/.agents/skills/akeyless-watchover
\end{Verbatim}
\exitcode{0}
\end{terminal}



\section{\texttt{watchover validate-skill}}\label{kit:validate-skill}
Check a skill folder against the Agent Skills specification. \emph{Offline.}

\begin{Verbatim}[fontsize=\footnotesize,breaklines,frame=leftline,framerule=1pt,rulecolor=\color{wo-rule}]
Validate a skill folder against the Agent Skills specification (frontmatter fields and limits, name = folder, body length, referenced files).
\end{Verbatim}

\paragraph{Options}\leavevmode
{\small
\begin{longtable}{@{}>{\raggedright\arraybackslash}p{0.27\linewidth}>{\raggedright\arraybackslash}p{0.17\linewidth}>{\raggedright\arraybackslash}p{0.13\linewidth}>{\raggedright\arraybackslash}p{0.35\linewidth}@{}}
\toprule
\textbf{Option} & \textbf{Value} & \textbf{Required / default} & \textbf{Description} \\ \midrule\endhead
\texttt{path} & \texttt{PATH} & \textbf{required} & skill folder (default: this skill) \\
\bottomrule
\end{longtable}}

\paragraph{Usage patterns}\leavevmode
\begin{Verbatim}[fontsize=\footnotesize,breaklines,breakanywhere]
scripts/watchover validate-skill
scripts/watchover validate-skill dist/akeyless-watchover
\end{Verbatim}
\subsection*{Example execution}\addcontentsline{toc}{subsection}{Example execution}
% GENERATED by docs/manual/build.py - do not edit; change the source it names instead.
% source: docs/manual/build.py EXAMPLES (executed in a sandbox copy of the repository)
The skill source:
\begin{terminal}
\begin{Verbatim}[fontsize=\scriptsize,breaklines,breakanywhere,commandchars=\\\{\}]
\textcolor{wo-prompt}{\$} \textbf{scripts/watchover validate-skill}
ok    skill/akeyless-watchover follows the Agent Skills specification
\end{Verbatim}
\exitcode{0}
\end{terminal}

A copy with an invalid name (exit code 1): (The copy is created before and removed after the command shown.)
\begin{terminal}
\begin{Verbatim}[fontsize=\scriptsize,breaklines,breakanywhere,commandchars=\\\{\}]
\textcolor{wo-prompt}{\$} \textbf{scripts/watchover validate-skill /tmp/wo-broken      # name: Akeyless_Watchover, folder wo-broken}
FAIL  name: 1-64 characters, lowercase letters, digits and single hyphens, no leading/trailing hyphen
\end{Verbatim}
\exitcode{1}
\end{terminal}



\section{\texttt{watchover mcp}}\label{kit:mcp}
Serve the kit, the docs and the skill commands to any MCP client (stdio). \emph{Offline.}

\begin{Verbatim}[fontsize=\footnotesize,breaklines,frame=leftline,framerule=1pt,rulecolor=\color{wo-rule}]
MCP server for akeyless-watchover: the kit, the documentation and the skill commands over stdio.

Any MCP client (Claude Desktop/Code, OpenAI Codex, GitHub Copilot / VS Code, Cursor, Gemini CLI and
others) can use akeyless-watchover through it, whichever model runs the client:

  tools      watchover_<command> for every kit command (arguments generated from the command's options);
             commands that call Akeyless are only listed when WATCHOVER_MCP_LIVE=1
  resources  watchover://skill/SKILL.md (the skill instructions), watchover://docs/<file>.md,
             watchover://model/<file>.yaml (the YAML model of the kit root)
  prompts    one per skill command (design, change, review, ...): the skill instructions + your request

No third-party dependency: JSON-RPC 2.0 over stdin/stdout, one message per line (MCP stdio transport).
Configure the client to start:  python3 <repo>/scripts/watchover_mcp.py   (or: scripts/watchover mcp)
\end{Verbatim}

\paragraph{Options}\leavevmode
{\small
\begin{longtable}{@{}>{\raggedright\arraybackslash}p{0.27\linewidth}>{\raggedright\arraybackslash}p{0.17\linewidth}>{\raggedright\arraybackslash}p{0.13\linewidth}>{\raggedright\arraybackslash}p{0.35\linewidth}@{}}
\toprule
\textbf{Option} & \textbf{Value} & \textbf{Required / default} & \textbf{Description} \\ \midrule\endhead
\texttt{-\allowbreak{}{}-\allowbreak{}repo} & \texttt{REPO} & optional & kit root the tools run in and the model/docs are read from (default: \$WATCHOVER\_REPO or this repository) \\
\texttt{-\allowbreak{}{}-\allowbreak{}live} & -- & optional & also expose the commands that call Akeyless (same as WATCHOVER\_MCP\_LIVE=1) \\
\texttt{-\allowbreak{}{}-\allowbreak{}list} & -- & optional & print the tools, resources and prompts, then exit \\
\bottomrule
\end{longtable}}

\paragraph{Environment variables}\leavevmode\\
\begin{tabularx}{\linewidth}{@{}l>{\raggedright\arraybackslash}X@{}}
\toprule
\texttt{WATCHOVER\_\allowbreak{}MCP\_\allowbreak{}LIVE} & 1 = also expose the commands that call Akeyless (verify, export, account-settings, token) \\
\texttt{WATCHOVER\_\allowbreak{}REPO} & kit root the tools run in (default: this repository) \\
\bottomrule
\end{tabularx}

\paragraph{Usage patterns}\leavevmode
\begin{Verbatim}[fontsize=\footnotesize,breaklines,breakanywhere]
scripts/watchover mcp
WATCHOVER_MCP_LIVE=1 scripts/watchover mcp --repo ~/src/akeyless-governance
\end{Verbatim}
\subsection*{Example execution}\addcontentsline{toc}{subsection}{Example execution}
% GENERATED by docs/manual/build.py - do not edit; change the source it names instead.
% source: docs/manual/build.py EXAMPLES (executed in a sandbox copy of the repository)
What an MCP client will see (offline commands only):
\begin{terminal}
\begin{Verbatim}[fontsize=\scriptsize,breaklines,breakanywhere,commandchars=\\\{\}]
\textcolor{wo-prompt}{\$} \textbf{scripts/watchover mcp --list}
tools:     watchover_help, watchover_test, watchover_simulate, watchover_forwarder, watchover_conjur, watchover_build_skill, watchover_validate_skill
resources: watchover://skill/SKILL.md, watchover://docs/README.md, watchover://docs/agents.md, watchover://docs/api-vs-cli.md, watchover://docs/architecture.md, watchover://docs/brownfield-import.md, watchover://docs/cli.md, watchover://docs/conjur-migration.md, watchover://docs/limitations.md, watchover://docs/model-reference.md, watchover://docs/path-and-role-design.md, watchover://docs/pipeline.md, watchover://docs/pki.md, watchover://docs/rbac-model.md, watchover://docs/terraform.md, watchover://docs/testing-and-verification.md, watchover://docs/usage-guide.md, watchover://model/access_expectations.yaml, watchover://model/auth_methods.yaml, watchover://model/folders.yaml, watchover://model/gateway.yaml, watchover://model/governance.yaml, watchover://model/pki.yaml, watchover://model/roles/payments.yaml, watchover://model/roles/platform.yaml
prompts:   watchover-help, watchover-design, watchover-change, watchover-review, watchover-pipeline, watchover-import, watchover-conjur, watchover-pki, watchover-harden, watchover-gateway, watchover-verify, watchover-cli, watchover-explain
\end{Verbatim}
\exitcode{0}
\end{terminal}

With the live commands enabled:
\begin{terminal}
\begin{Verbatim}[fontsize=\scriptsize,breaklines,breakanywhere,commandchars=\\\{\}]
\textcolor{wo-prompt}{\$} \textbf{scripts/watchover mcp --list --live | head -n 1}
tools:     watchover_help, watchover_test, watchover_simulate, watchover_verify, watchover_export, watchover_forwarder, watchover_conjur, watchover_account_settings, watchover_token, watchover_build_skill, watchover_validate_skill
\end{Verbatim}
\exitcode{0}
\end{terminal}

A raw MCP session over stdio (what a client sends), output shortened per message: (mcp\_demo.py is the short Python script embedded in docs/manual/build.py.)
\begin{terminal}
\begin{Verbatim}[fontsize=\scriptsize,breaklines,breakanywhere,commandchars=\\\{\}]
\textcolor{wo-prompt}{\$} \textbf{python3 mcp_demo.py   # sends initialize, tools/call watchover_simulate and prompts/get watchover-change to `scripts/watchover mcp` over stdio and prints a summary of each answer}
initialize -> \{'name': 'watchover-mcp', 'title': 'akeyless-watchover', 'version': '1.5.0'\} protocol 2025-06-18
tools/call watchover_simulate -> isError False
  $ scripts/watchover simulate --model terraform/model --expectations terraform/model/access_expectations.yaml
  (exit code 0)
  ... 0 mismatch(es)
prompts/get watchover-change -> 19264 characters:
  Follow the akeyless-watchover skill instructions below. Run its `change` command with: give fraud-analysts rea...
\end{Verbatim}
\exitcode{0}
\end{terminal}



\section{\texttt{watchover help}}\label{kit:help}
Help in the terminal. \texttt{-h}/\texttt{--help} after any command shows that command; \texttt{watchover} with no argument is the same as \texttt{watchover help}.

\paragraph{Options}\leavevmode
{\small
\begin{longtable}{@{}>{\raggedright\arraybackslash}p{0.27\linewidth}>{\raggedright\arraybackslash}p{0.17\linewidth}>{\raggedright\arraybackslash}p{0.13\linewidth}>{\raggedright\arraybackslash}p{0.35\linewidth}@{}}
\toprule
\textbf{Option} & \textbf{Value} & \textbf{Required / default} & \textbf{Description} \\ \midrule\endhead
\texttt{topic} & \texttt{TOPIC} & \textbf{required} & a kit command name or 'skill'; omitted = overview \\
\texttt{-\allowbreak{}{}-\allowbreak{}markdown} & -- & optional & print the whole usage guide as Markdown (docs/usage-guide.md) \\
\bottomrule
\end{longtable}}

\subsection*{Example execution}\addcontentsline{toc}{subsection}{Example execution}
% GENERATED by docs/manual/build.py - do not edit; change the source it names instead.
% source: docs/manual/build.py EXAMPLES (executed in a sandbox copy of the repository)
\begin{terminal}
\begin{Verbatim}[fontsize=\scriptsize,breaklines,breakanywhere,commandchars=\\\{\}]
\textcolor{wo-prompt}{\$} \textbf{scripts/watchover help}
watchover 1.5.1 - akeyless-watchover kit

Usage: scripts/watchover <command> [options]

Commands:
  test              Run every offline check (what CI runs).
  simulate          Offline: evaluate access_expectations.yaml against the model.
  verify            Live: compare expectations with describe-permissions run as each persona's CLI profile.  [needs Akeyless]
  export            Export existing config and translate it: YAML model + import blocks, or plain Terraform (--hcl).  [needs Akeyless]
  forwarder         Scaffold a new event forwarder (model entry or HCL); --import-existing adopts one by name.
  conjur            Translate Conjur policy into model roles + identity bindings.
  account-settings  Apply account_settings_cli (settings the provider cannot manage) with the CLI.  [needs Akeyless]
[... 12 lines omitted ...]
\end{Verbatim}
\exitcode{0}
\end{terminal}



\section{\texttt{watchover version}}\label{kit:version}
Prints the kit version (also \texttt{--version}). No options.

\subsection*{Example execution}\addcontentsline{toc}{subsection}{Example execution}
% GENERATED by docs/manual/build.py - do not edit; change the source it names instead.
% source: docs/manual/build.py EXAMPLES (executed in a sandbox copy of the repository)
\begin{terminal}
\begin{Verbatim}[fontsize=\scriptsize,breaklines,breakanywhere,commandchars=\\\{\}]
\textcolor{wo-prompt}{\$} \textbf{scripts/watchover version}
1.5.0
\end{Verbatim}
\exitcode{0}
\end{terminal}



